%! Unit 2, Lessons 2.1–2.5 — ACTUAL Unit Test (Answer Key)
% The scored unit test. Item for item it is the parallel form of ../practice_test/: same two
% sections, the same ten multiple-choice targets in the same order, the same two free-response
% questions with the same parts and point values — every context different, so a student who
% worked the practice test has practised the moves and not memorised the answers.
% BLUEPRINT (1 pt per MC, FRQ 1 = 4, FRQ 2 = 6, total 20; one 55-minute period):
%   MC 1 conditional relative frequency (denominator)  — clinic walk-in vs scheduled × wait    2.1
%   MC 2 association, groups of unequal size           — street vs park trees × stressed       2.1
%   MC 3 explanatory vs response                       — library hours open vs checkouts       2.2
%   MC 4 describe a scatterplot                        — river water temp vs dissolved oxygen     2.2
%   MC 5 properties of r                               — hardware store floor area vs sales    2.2
%   MC 6 predict + extrapolation                       — tomato plant age vs height             2.3
%   MC 7 residual, sign and meaning                    — dairy feed vs milk                    2.3
%   MC 8 slope in context                              — lake algae vs water clarity           2.4
%   MC 9 r^2 from computer output                      — home size vs monthly heating cost     2.4
%   MC 10 residual plot vs high r                      — fish-farm age vs weight               2.5
%   FRQ 1 two-way table                                — rail-station distance × commute mode  2.1
%   FRQ 2 regression                                   — growing-season rainfall vs wheat yield 2.2–2.5
% THE DATA
%   MC 1: walk-in 54/36 (90), scheduled 42/168 (210), totals 96/204/300 -> 54/90 = 0.60.
%   MC 9: n = 30, xbar 2000, s_x 500, ybar 150, s_y 40, r 0.77 -> b 0.0616, a 26.80, S 25.97,
%         SE(b) 0.009646, T 6.39; SE(a) 19.87, T 1.35, P 0.188; R-Sq 59.3%, adj 57.8%.
%   FRQ 1: near 96/64/40 (200), far 54/216/30 (300); transit 150, car 280, bike/walk 70; 500.
%   FRQ 2: 15 fields, rainfall x (in): 9 10 11 12 13 14 14 15 16 17 18 19 20 21 22
%          yield y (bu/acre):          43 46 44 49 53 58 53 55 62 57 60 66 68 64 71
%          yhat = 24.85 + 2.062x, r = 0.957, R-Sq 91.6%, S 2.624, SE(b) 0.1732, T 11.91,
%          SE(a) 2.752, T 9.03. Residuals: -.40 .53 -3.53 -.59 1.35 4.29 -.71 -.78 4.16 -2.90
%          -1.96 1.98 1.92 -4.15 .79 (random scatter). The (14, 58) residual is 4.28 with the
%          printed coefficients.
% NAMESTRIP: \namedateperiod mirrors the blank — structural, so the pages stay in lockstep.
% NO teachernote here — scoring guidance belongs in the teacher's copy, never in a key.
% GENERATED FROM the blank.
% PAGINATION: each item carries a \boxguard sized to its height so no stem, graph or option
% list splits across a page; the Section II strip is guarded to travel with FRQ 1.
% NO SKETCH ITEMS: every display is pre-drawn in TikZ and read.
% PAGE LOCKSTEP with ../../tests/actual_test/: every \writelines{n} here is answered by
% exactly n \ansline{} there, each terse enough not to wrap; every work block is byte-identical.
% Nothing else differs except the preamble, the header, and the \ans{} marks on the correct MC
% options.
\documentclass[12pt]{article}
\usepackage{apstats-article}
\usepackage{apstats-key}   % pulls in -boxes; do NOT also load -boxes

\begin{document}

\pageheader{Unit 2 \quad Lessons 2.1--2.5}{Unit Test --- Answer Key}

\namedateperiod
\par\vspace{2pt}
\hfill\small Score:\;\blank{1.3cm}\;/\;20\normalsize
\vspace{0.06in}

\boxguard[10]
\begin{headlinebox}{goldbg}
\small
\textbf{Section I --- Multiple Choice} (1 point each). Circle the single best answer. Read all
five options --- more than one of them is about the context, not just the arithmetic.
\end{headlinebox}

\vspace{0.10in}

\begin{enumerate}[leftmargin=1.9em, itemsep=0.13in]

  \item A health clinic recorded, for 300 patient visits last month, the type of visit and
        whether the patient waited more than 30 minutes to be seen.

        \vspace{2pt}
        {\centering\small
        \renewcommand{\arraystretch}{1.15}
        \begin{tabular}{@{} l c c c @{}}
        \rowcolor{sky}
         & \textbf{Waited over 30 min} & \textbf{Did not} & \textbf{Total} \\
        \midrule
        Walk-in   & 54 & 36  & 90  \\
        Scheduled & 42 & 168 & 210 \\
        \midrule
        Total     & 96 & 204 & 300 \\
        \end{tabular}\par}
        \vspace{4pt}

        What proportion of \textbf{walk-in} visits waited more than 30 minutes?
        \par\vspace{3pt}\noindent
        \hspace{1.2em}\textbf{(A)} 0.18 \hfill \textbf{(B)} 0.30 \hfill
        \textbf{(C)} 0.32 \hfill \textbf{(D)} 0.5625 \hfill \textbf{(E)} \ans{0.60}

  \boxguard[18]% keep the item whole
  \item A city forester checked 240 street trees and 60 park trees and rated each one
        \emph{stressed} or \emph{healthy}. The segmented bar graph shows the results.

        \vspace{4pt}
        {\centering
        \begin{tikzpicture}[x=1cm, y=0.024cm]
          \foreach \y in {0,50,100} {
            \draw (-0.08,\y) -- (0.08,\y) node[left, font=\tiny, xshift=-2pt] {\y\%};}
          \draw[thick] (0,0) -- (0,104);
          \draw[thick] (0,0) -- (3.6,0);
          \fill[navy]  (0.4,0)  rectangle (1.6,35);
          \fill[sky]   (0.4,35) rectangle (1.6,100);
          \fill[navy]  (2.1,0)  rectangle (3.3,15);
          \fill[sky]   (2.1,15) rectangle (3.3,100);
          \draw (0.4,0) rectangle (1.6,100);
          \draw (2.1,0) rectangle (3.3,100);
          \node[white, font=\scriptsize\bfseries] at (1.0,17.5) {35\%};
          \node[font=\scriptsize\bfseries] at (1.0,67.5) {65\%};
          \node[white, font=\tiny\bfseries] at (2.7,7.5) {15\%};
          \node[font=\scriptsize\bfseries] at (2.7,57.5) {85\%};
          \node[below, font=\tiny, align=center] at (1.0,-2) {Street trees\\($n = 240$)};
          \node[below, font=\tiny, align=center] at (2.7,-2) {Park trees\\($n = 60$)};
          \fill[navy] (4.0,80) rectangle (4.25,90); \node[right, font=\tiny] at (4.3,85) {Stressed};
          \fill[sky]  (4.0,62) rectangle (4.25,72); \draw (4.0,62) rectangle (4.25,72);
          \node[right, font=\tiny] at (4.3,67) {Healthy};
        \end{tikzpicture}\par}
        \vspace{2pt}

        Which statement is best supported by these data?
        \begin{enumerate}[label=\textbf{(\Alph*)}, leftmargin=2.2em, itemsep=2pt]
          \item Street trees are more often stressed, because 84 street trees are stressed and
                only 9 park trees are.
          \item No comparison is possible, because there are four times as many street trees.
          \item \ans{Location and tree health are associated: 35\% of street trees are stressed,
                compared with 15\% of park trees.}
          \item Location and tree health are not associated, because most trees in both
                locations are healthy.
          \item Growing along a street causes 35\% of trees to become stressed.
        \end{enumerate}

  \boxguard[9]% keep the item whole
  \item A library system records, for each of its 18 branches, the number of hours the branch
        is open per week and the number of books checked out per week. The staff want to know
        whether longer hours go with more checkouts. Which choice of variables is correct?
        \begin{enumerate}[label=\textbf{(\Alph*)}, leftmargin=2.2em, itemsep=2pt]
          \item \ans{Explanatory: hours open per week. Response: books checked out per week.}
          \item Explanatory: books checked out per week. Response: hours open per week.
          \item Explanatory: the 18 branches. Response: books checked out per week.
          \item Both are response variables, because both are recorded every week.
          \item Neither can be named until the correlation has been computed.
        \end{enumerate}

  \boxguard[18]% keep the item whole
  \item The scatterplot shows the water temperature and the dissolved oxygen measured at 15
        river sites on one summer afternoon.

        \vspace{4pt}
        {\centering
        \begin{tikzpicture}[xscale=0.20, yscale=0.28]
          \foreach \y in {6,8,10,12} {\draw[linegray, very thin] (0,\y) -- (29,\y);
            \node[left, font=\tiny] at (0,\y) {\y};}
          \draw[->, thick] (0,4.5) -- (30,4.5);
          \draw[->, thick] (0,4.5) -- (0,14);
          \foreach \x in {0,5,10,15,20,25} {\draw (\x,4.3) -- (\x,4.7);
            \node[below, font=\tiny] at (\x,4.4) {\x};}
          \foreach \x/\y in {4/12.6, 6/12.1, 8/11.9, 10/11.2, 11/11.0, 13/10.4, 14/10.3, 16/9.7,
                             17/9.6, 19/9.0, 20/8.9, 22/8.3, 23/8.2, 25/7.7}
            {\fill[navy] (\x,\y) circle[x radius=0.32, y radius=0.23];}
          \fill[navy] (9,6.0) circle[x radius=0.32, y radius=0.23];
          \node[font=\tiny] at (14.5,2.2) {Water temperature ($^\circ$C)};
          \node[rotate=90, font=\tiny] at (-4.3,9.2) {Dissolved oxygen (mg/L)};
        \end{tikzpicture}\par}
        \vspace{2pt}

        Which is the best description of the association?
        \begin{enumerate}[label=\textbf{(\Alph*)}, leftmargin=2.2em, itemsep=2pt]
          \item Strong, positive, and linear, with no unusual points.
          \item Strong, negative, and linear, with no unusual points.
          \item \ans{Strong, negative, and linear, with one site far below the pattern.}
          \item Weak and curved, because one site does not follow the others.
          \item There is no association, because one site has low temperature and low
                oxygen.
        \end{enumerate}

  \boxguard[9]% keep the item whole
  \item A hardware chain finds $r = 0.74$ between store floor area (square feet) and monthly
        sales (dollars) for its 40 stores. Which statement is true?
        \begin{enumerate}[label=\textbf{(\Alph*)}, leftmargin=2.2em, itemsep=2pt]
          \item If floor area were recorded in square meters, $r$ would change.
          \item \ans{If sales were in thousands of dollars and floor area in square meters, $r$
                would still be 0.74.}
          \item Using sales to predict floor area would give $r = 1/0.74 \approx 1.35$.
          \item $r = 0.74$ means that 74\% of the stores have sales that match their size.
          \item $r$ is measured in dollars per square foot.
        \end{enumerate}

  \boxguard[10]% keep the item whole
  \item A gardener fits $\hat{y} = 2.5 + 1.2x$, where $x$ is a tomato plant's age (days) and
        $\hat{y}$ is its predicted height (cm), using plants 10 to 40 days old. What does the
        line say about a plant 120 days old? Which statement is best?
        \begin{enumerate}[label=\textbf{(\Alph*)}, leftmargin=2.2em, itemsep=2pt]
          \item 146.5 cm, and the prediction is reliable because it comes from the line.
          \item No number can be computed, because no plant in the data was 120 days old.
          \item 97.9 days, found by solving $120 = 2.5 + 1.2x$.
          \item 146.5 cm, and every 120-day-old plant will be exactly that tall.
          \item \ans{146.5 cm, but this is extrapolation --- 120 days is far outside the 10 to 40
                days in the data --- so the prediction may be very unreliable.}
        \end{enumerate}

  \boxguard[9]% keep the item whole
  \item On a dairy farm, $\hat{y} = 4.2 + 1.5x$ predicts a cow's daily milk (liters) from its
        daily feed (kilograms). A cow that ate 20 kg of feed gave 31 liters of milk. What is
        its residual, and what does it mean?
        \begin{enumerate}[label=\textbf{(\Alph*)}, leftmargin=2.2em, itemsep=2pt]
          \item 3.2 liters; the cow gave more milk than predicted.
          \item $-3.2$ liters; the cow gave more milk than predicted.
          \item 3.2 liters; the cow gave less milk than predicted.
          \item \ans{$-3.2$ liters; the cow gave less milk than predicted.}
          \item 34.2 liters; the line predicted this cow exactly.
        \end{enumerate}

  \boxguard[11]% keep the item whole
  \item For 25 lakes, $\hat{y} = 18.4 - 0.62x$, where $x$ is algae concentration (micrograms
        per liter) and $\hat{y}$ is predicted water clarity (feet). Which is the correct
        interpretation of the slope?
        \begin{enumerate}[label=\textbf{(\Alph*)}, leftmargin=2.2em, itemsep=2pt]
          \item \ans{For each additional microgram per liter of algae, the predicted water clarity
                decreases by 0.62 feet.}
          \item For each additional microgram per liter of algae, a lake's water clarity
                decreases by exactly 0.62 feet.
          \item For each additional foot of clarity, the predicted algae concentration
                decreases by 0.62 micrograms per liter.
          \item When a lake has no algae, its water clarity is 18.4 feet.
          \item Adding algae to a lake causes its clarity to fall 0.62 feet per microgram.
        \end{enumerate}

  \boxguard[15]% keep the item whole
  \item An energy company regressed monthly winter heating cost (dollars) on home size
        (square feet) for 30 homes.

        \vspace{2pt}
        {\centering\footnotesize\ttfamily
        \begin{tabular}{@{} l r r r r @{}}
        Predictor & Coef & SE Coef & T & P \\ \hline
        Constant & 26.80 & 19.87 & 1.35 & 0.188 \\
        Size & 0.06160 & 0.009646 & 6.39 & 0.000 \\
        \end{tabular}\par
        S = 25.97 \quad R-Sq = 59.3\% \quad R-Sq(adj) = 57.8\%\par}
        \vspace{2pt}

        Which is the correct interpretation of the value of R-Sq?
        \begin{enumerate}[label=\textbf{(\Alph*)}, leftmargin=2.2em, itemsep=2pt]
          \item The correlation between home size and heating cost is 0.593.
          \item About 59.3\% of the homes have heating costs the line predicts correctly.
          \item \ans{About 59.3\% of the variation in monthly heating cost is explained by the
                linear model with home size.}
          \item Home size causes 59.3\% of a home's heating cost.
          \item Each extra square foot raises the predicted heating cost by 59.3\%.
        \end{enumerate}

  \boxguard[18]% keep the item whole
  \item A fish farm fits a line predicting a fish's weight (grams) from its age (weeks) for
        14 fish, and finds $r = 0.96$. The residual plot is below.

        \vspace{4pt}
        {\centering
        \begin{tikzpicture}[xscale=0.19, yscale=0.055]
          \draw[->, thick] (0,-24) -- (0,28);
          \draw[thick] (0,-24) -- (30,-24);
          \draw[dashed, slate] (0,0) -- (30,0);
          \foreach \y in {-20,0,20} {\draw (-0.4,\y) -- (0.4,\y);
            \node[left, font=\tiny] at (0,\y) {\y};}
          \foreach \x in {0,10,20,30} {\draw (\x,-25) -- (\x,-23);
            \node[below, font=\tiny] at (\x,-24) {\x};}
          \foreach \x/\e in {2/24, 4/14, 6/5, 8/-2, 10/-8, 12/-13, 14/-16, 16/-16, 18/-14,
                             20/-10, 22/-4, 24/4, 26/14, 28/24}
            {\fill[navy] (\x,\e) circle[x radius=0.35, y radius=1.2];}
          \node[font=\tiny] at (15,-34) {Age (weeks)};
          \node[rotate=90, font=\tiny] at (-5.2,2) {Residual (g)};
        \end{tikzpicture}\par}
        \vspace{2pt}

        Which conclusion is best?
        \begin{enumerate}[label=\textbf{(\Alph*)}, leftmargin=2.2em, itemsep=2pt]
          \item A line is appropriate, because $r = 0.96$ is close to 1.
          \item Age and weight are not associated, because the residuals are not all zero.
          \item A line is appropriate, because there are about as many positive residuals as
                negative ones.
          \item \ans{A line is not appropriate: the residuals form a clear curve, even though $r$
                is high.}
          \item The curve in the residuals shows that the farm weighed the fish incorrectly.
        \end{enumerate}

\end{enumerate}

\vspace{0.14in}

\boxguard[26]
\begin{headlinebox}{goldbg}
\small
\textbf{Section II --- Free Response} (10 points). Show your reasoning. Every answer must be
written \emph{in the context of the problem} --- name the variables and their units.
\end{headlinebox}

\vspace{0.10in}

\begin{enumerate}[leftmargin=1.9em, itemsep=0.16in, resume]

  \boxguard[18]% keep the item whole
  \item \textbf{(4 points)} A regional transit agency surveyed 500 commuters about whether they
        live within half a mile of a rail station and how they usually get to work.

        \vspace{4pt}
        \noindent
        \begin{minipage}[c]{0.56\linewidth}\centering\small
        \renewcommand{\arraystretch}{1.2}
        \setlength{\tabcolsep}{4pt}
        \begin{tabular}{@{} l c c c c @{}}
        \rowcolor{sky}
         & \textbf{Transit} & \textbf{Car} & \textbf{Bike/walk} & \textbf{Total} \\
        \midrule
        Near a station & 96 & 64  & 40 & 200 \\
        Far from one   & 54 & 216 & 30 & 300 \\
        \midrule
        Total          & 150 & 280 & 70 & 500 \\
        \end{tabular}
        \end{minipage}\hfill
        \begin{minipage}[c]{0.42\linewidth}\centering
        \begin{tikzpicture}[x=1cm, y=0.032cm]
          \foreach \y in {0,50,100} {
            \draw (-0.08,\y) -- (0.08,\y) node[left, font=\tiny, xshift=-2pt] {\y\%};}
          \draw[thick] (0,0) -- (0,104);
          \draw[thick] (0,0) -- (2.9,0);
          \fill[navy]    (0.3,0)  rectangle (1.3,48);
          \fill[goldacc] (0.3,48) rectangle (1.3,80);
          \fill[sky]     (0.3,80) rectangle (1.3,100);
          \fill[navy]    (1.7,0)  rectangle (2.7,18);
          \fill[goldacc] (1.7,18) rectangle (2.7,90);
          \fill[sky]     (1.7,90) rectangle (2.7,100);
          \draw (0.3,0) rectangle (1.3,100);
          \draw (1.7,0) rectangle (2.7,100);
          \node[white, font=\tiny\bfseries] at (0.8,24) {48\%};
          \node[font=\tiny\bfseries] at (0.8,64) {32\%};
          \node[font=\tiny\bfseries] at (0.8,90) {20\%};
          \node[white, font=\tiny\bfseries] at (2.2,9) {18\%};
          \node[font=\tiny\bfseries] at (2.2,54) {72\%};
          \node[font=\tiny\bfseries] at (2.2,95) {10\%};
          \node[below, font=\tiny] at (0.8,-2) {Near};
          \node[below, font=\tiny] at (2.2,-2) {Far};
          \fill[sky]     (3.0,86) rectangle (3.22,96); \draw (3.0,86) rectangle (3.22,96);
          \node[right, font=\tiny] at (3.25,91) {Bike/walk};
          \fill[goldacc] (3.0,70) rectangle (3.22,80); \node[right, font=\tiny] at (3.25,75) {Car};
          \fill[navy]    (3.0,54) rectangle (3.22,64); \node[right, font=\tiny] at (3.25,59) {Transit};
        \end{tikzpicture}
        \end{minipage}

        \vspace{6pt}
        \textbf{(a)} What proportion of \emph{all} 500 commuters use transit? Of the commuters
        who live \emph{far} from a station, what proportion drive a car?
        \begin{work}
          \text{Transit, all commuters: } &\ 150/500 = 0.30 \\
          \text{Car, among far commuters: } &\ 216/300 = 0.72
        \end{work}

        \vspace{2pt}
        \boxguard[6]% keep the item whole
        \textbf{(b)} Use the segmented bar graph to compare the distributions of commute mode
        for commuters who live near a station and those who live far from one.\\[4pt]
        \ansline{Near: 48\% transit, 32\% car, 20\% bike/walk. Far: 18\% transit, 72\% car,}\\
        \ansline{10\% bike/walk. Near commuters take transit far more often (48\% vs.\ 18\%);}\\
        \ansline{far commuters mostly drive (72\% vs.\ 32\%) and bike or walk less (10\% vs.\ 20\%).}

        \vspace{6pt}
        \boxguard[6]% keep the item whole
        \textbf{(c)} Is there an association between living near a rail station and commute
        mode for these 500 commuters? Justify with numbers. Can the agency conclude that
        living near a station \emph{causes} people to take transit? Explain.\\[4pt]
        \ansline{Yes: the mode distributions differ a lot --- 48\% of near commuters take transit,}\\
        \ansline{only 18\% of far ones. No cause: this is a survey, not an experiment; people who}\\
        \ansline{like transit may choose to live near a station, or other variables may differ.}

  \boxguard[42]% keep the item whole
  \item \textbf{(6 points)} An agricultural extension agent recorded the growing-season rainfall
        (inches) and the wheat yield (bushels per acre) for 15 fields in one county. The
        scatterplot, computer output, and residual plot are below.

        \vspace{4pt}
        \noindent
        \begin{minipage}[b]{0.49\linewidth}\centering
        \begin{tikzpicture}[xscale=0.38, yscale=0.075]
          \foreach \y in {40,50,60,70} {\draw[linegray, very thin] (8,\y) -- (23,\y);
            \node[left, font=\tiny] at (8,\y) {\y};}
          \draw[->, thick] (8,37) -- (23.5,37);
          \draw[->, thick] (8,37) -- (8,75);
          \foreach \x in {10,14,18,22} {\draw (\x,36.3) -- (\x,37.7);
            \node[below, font=\tiny] at (\x,36.6) {\x};}
          \foreach \x/\y in {9/43, 10/46, 11/44, 12/49, 13/53, 14/58, 14/53, 15/55, 16/62,
                             17/57, 18/60, 19/66, 20/68, 21/64, 22/71}
            {\fill[navy] (\x,\y) circle[x radius=0.17, y radius=0.85];}
          \node[font=\tiny] at (15.5,29.5) {Rainfall (inches)};
          \node[rotate=90, font=\tiny] at (5.6,56) {Yield (bu/acre)};
        \end{tikzpicture}
        \end{minipage}\hfill
        \begin{minipage}[b]{0.49\linewidth}\centering
        \begin{tikzpicture}[xscale=0.38, yscale=0.25]
          \draw[->, thick] (8,-5.5) -- (8,6);
          \draw[thick] (8,-5.5) -- (23.5,-5.5);
          \draw[dashed, slate] (8,0) -- (23.5,0);
          \foreach \y in {-4,-2,0,2,4} {\draw (7.8,\y) -- (8.2,\y);
            \node[left, font=\tiny] at (8,\y) {\y};}
          \foreach \x in {10,14,18,22} {\draw (\x,-5.65) -- (\x,-5.35);
            \node[below, font=\tiny] at (\x,-5.5) {\x};}
          \foreach \x/\e in {9/-0.40, 10/0.53, 11/-3.53, 12/-0.59, 13/1.35, 14/4.29, 14/-0.71,
                             15/-0.78, 16/4.16, 17/-2.90, 18/-1.96, 19/1.98, 20/1.92, 21/-4.15,
                             22/0.79}
            {\fill[navy] (\x,\e) circle[x radius=0.17, y radius=0.26];}
          \node[font=\tiny] at (15.5,-7.8) {Rainfall (inches)};
          \node[rotate=90, font=\tiny] at (5.9,0) {Residual (bu/acre)};
        \end{tikzpicture}
        \end{minipage}

        \vspace{4pt}
        {\centering\footnotesize\ttfamily
        \begin{tabular}{@{} l r r r r @{}}
        Predictor & Coef & SE Coef & T & P \\ \hline
        Constant & 24.85 & 2.752 & 9.03 & 0.000 \\
        Rainfall & 2.062 & 0.1732 & 11.91 & 0.000 \\
        \end{tabular}\par
        S = 2.624 \quad R-Sq = 91.6\% \quad R-Sq(adj) = 91.0\%\par}

        \vspace{6pt}
        \textbf{(a)} Describe the association between rainfall and wheat yield.\\[4pt]
        \ansline{Strong, positive, linear association ($r = \sqrt{0.916} \approx 0.957$): fields with}\\
        \ansline{more rainfall tend to have higher wheat yields. No clear outliers.}

        \vspace{6pt}
        \textbf{(b)} Write the equation of the least-squares regression line, using the
        variable names. Then interpret its slope in context.\\[4pt]
        \ansline{$\widehat{\text{yield}} = 24.85 + 2.062(\text{rainfall})$}\\
        \ansline{For each additional inch of growing-season rainfall, the \emph{predicted} wheat}\\
        \ansline{yield increases by 2.062 bushels per acre.}

        \vspace{6pt}
        \textbf{(c)} One field had 14 inches of rainfall and yielded 58 bushels per acre. Find
        its residual and interpret it in context.
        \begin{work}
          \widehat{\text{yield}} &= 24.85 + 2.062(14) = 53.72 \text{ bu/acre} \\
          \text{residual} &= y - \hat{y} = 58 - 53.72 = 4.28 \text{ bu/acre}
        \end{work}
        \ansline{This field yielded 4.28 bu/acre \emph{more} than predicted for 14 inches of rain.}

        \vspace{6pt}
        \textbf{(d)} Interpret the value of R-Sq in context.\\[4pt]
        \ansline{About 91.6\% of the variation in wheat yield among these 15 fields is explained}\\
        \ansline{by the linear model with growing-season rainfall.}

        \vspace{6pt}
        \textbf{(e)} Use the residual plot to decide whether a linear model is appropriate for
        these data. Explain.\\[4pt]
        \ansline{Yes: the residuals scatter randomly above and below 0 with no curve and no}\\
        \ansline{fanning, so a linear model for yield from rainfall is appropriate.}

\end{enumerate}

\end{document}
